% !TEX program = xelatex
% 问题三 完整实现方案（确定性两阶段机器狗搜索与清除）
% 编译：xelatex main.tex && xelatex main.tex
\documentclass[a4paper,11pt]{ctexart}

\usepackage[margin=2.4cm]{geometry}
\usepackage{amsmath,amssymb,bm}
\usepackage{booktabs}
\usepackage{graphicx}
\usepackage{float}
\usepackage{array}
\usepackage{multirow}
\usepackage{caption}
\usepackage{hyperref}
\usepackage{xcolor}
\usepackage{listings}
\usepackage{tikz}
\usetikzlibrary{arrows.meta,calc}

\captionsetup{font=small,labelfont=bf}
\hypersetup{colorlinks=true,linkcolor=blue!60!black,citecolor=blue!60!black,urlcolor=blue!60!black}

\newcommand{\R}{\mathbb{R}}
\newcommand{\code}[1]{\texttt{#1}}
\newcommand{\figref}[1]{图~\ref{#1}}
\newcommand{\tabref}[1]{表~\ref{#1}}
\newcommand{\eqrefn}[1]{式~(\ref{#1})}

\title{\bfseries 问题三 完整实现方案\\[4pt]
\large 机器狗搜索与清除干扰源的确定性两阶段策略}
\author{}
\date{}

\begin{document}
\maketitle
\vspace{-2.5em}
\begin{center}
\small 本文档描述问题三确定性求解方案的完整实现：覆盖圆布局、起始全频道扫描与布局定向、\\[2pt]
巡视扫描与顺路清除、阶段二定位与清除，以及验证与实测结果。全部流程零随机、同种子逐字节可复现。
\end{center}

\tableofcontents
\newpage

%======================================================================
\section{问题背景与任务}
%======================================================================
竞赛方在半径 $R=1800$\,m 的圆域内随机布置若干干扰源（频率通道 $1\sim 20$），每局 $10\sim 16$ 个源，
每个源有固定的频道、位置、类型（全向/定向）与接收半径 $\rho_{\text{rec}}\in[1000,1500]$\,m。
机器狗在圆域原点出发，通过以下接口交互：
\begin{itemize}
  \item \code{/measure}：在位置 $(x,y)$ 对频道 $c$ 测向，返回
        $\text{result}\in\{\texttt{no\_signal},\texttt{near},\texttt{direction}\}$；
        返回 \texttt{direction} 时给出示向度 $\theta$，误差 $|e|\le 1^\circ$，且\textbf{同一位置、
        同一频道重复测量误差完全相同}（空间相关系统误差，非独立随机误差）；
  \item \code{/clear}：在位置 $(x,y)$ 对频道 $c$ 光学精确定位，成功要求狗与源距离 $\le 20$\,m；
  \item 移动速度 $5$\,m/s；现实限时 $20$ 分钟/局（求解须考虑预留余量）。
\end{itemize}

\textbf{任务}：在干扰源真值不可见的前提下，规划狗的巡视与测量，把\textbf{全部}干扰源清除，
并尽量缩短虚拟总时长（移动 + 测向 + 清除的虚拟时钟总和）。

关键难点：
\begin{enumerate}
  \item 单次测向有沿视线方向的偏移（$|e|\le1^\circ$），单视角定位不可靠，必须\textbf{多视角交会}；
  \item 同点复测无信息增益，必须\textbf{换位置才能获取新信息}；
  \item 接收半径 1000–1500\,m 并未给出，源可能在 1000\,m 之外，存在大量"测了收不到"的频道；
  \item 源到每个观测站的距离未知，交会区域可能无界或病态，需要末端逼近机制兜底。
\end{enumerate}

%======================================================================
\section{总体方案架构}
%======================================================================
本方案采用\textbf{确定性两阶段}结构，全程无随机数，同一随机种子下逐字节可复现：

\begin{figure}[H]
  \centering
  \begin{tikzpicture}[
      box/.style={rectangle,draw,rounded corners=2pt,align=center,minimum height=9mm,fill=blue!6},
      arr/.style={-{Stealth[length=2.6mm]},thick},
    ]
    \node[box] (c0) at (0,0) {\textbf{覆盖圆布局}；\\[1pt] 7 站、$\rho_{\text{worst}}\le 1000$\,m};
    \node[box] (c1) at (5.2,0) {\textbf{起点全频道扫描}\\[1pt] 20 频道 $\approx 119$\,s};
    \node[box] (c2) at (10.4,0) {\textbf{布局定向}\\[1pt] 密集 $60^\circ$ 扇区 $\to$ 第一点 $\theta^*-90^\circ$};
    \node[box] (c3) at (15.5,0) {\textbf{阶段一 巡视扫描}\\[1pt] 7 站测向 + 顺路清除};
    \node[box] (c4) at (15.5,-2.6) {\textbf{阶段二 定位与清除}\\[1pt] 就近试清 $\to$ 补测 $\to$ 逼近};
    \node[box] (c5) at (0,-2.6) {\textbf{完成}\\[1pt] 全部源清除};
    \draw[arr] (c0) -- (c1);
    \draw[arr] (c1) -- (c2);
    \draw[arr] (c2) -- (c3);
    \draw[arr] (c3) -- (c4);
    \draw[arr] (c4) -- (c5);
    \draw[arr] (c2) to[bend left=25] (c4);
  \end{tikzpicture}
  \caption{方案总体流程。阶段一在巡视途中顺路清除可判定的源，把"专程跑一趟"变成"顺路清掉"。}
  \label{fig:flow}
\end{figure}

\begin{itemize}
  \item \textbf{阶段一（巡视扫描）}：按覆盖圆布局巡视 7 个观测站，每站对"仍不够准"的频道测向，
        采集多视角示向度；站点间移动时按方位扇区判据\textbf{顺路清除}可判定的源；
  \item \textbf{阶段二（定位与清除）}：对剩余每个频道执行"就近试清 $\to$ 多点覆盖试清 $\to$
        文献准则补测缩小区域 $\to$ 沿最新示向度逼近"的四级流程，直到清除成功。
\end{itemize}

两阶段之间的\textbf{衔接设计}：巡视终点（由最短开放路径 + 源密集方向共同决定）尽量靠近源密集
区域，使阶段二的前几个频道少跑路；布局旋转是零成本自由度（覆盖与路径长度都不依赖朝向），
用于把第一观测点对准源密集方向的旁侧。

%======================================================================
\section{覆盖圆布局}
%======================================================================
\subsection{布局求解}
利用问题一的覆盖思想：只要在圆域内布置若干观测站，使\textbf{任意点到最近观测站的距离} $\le
\rho_{\min}=1000$\,m（有效接收半径下界），则该源总能被至少一个站听到方向。经网格 + 边界 + 解析
候选点三重校验，采用 \textbf{7 站均匀布局}：7 个站等分在半径 $1000$\,m 的圆上（正七边形，
相邻站位间隔 $51.43^\circ$），最坏最近距离恰好 $1000.0$\,m、最坏点即区域圆心，覆盖率 $100\%$、
平均重数 $1.96$——这是"$\le\rho_{\min}$"的\textbf{临界紧解}（与 7 个半径 $1000$\,m 圆的面积
上界吻合，余量为 $0$）。作为对照，另有数值优化的 7 点一般布局（里程 $6167$\,m、余量 $\sim5$\,m、
站点距圆心 $975\sim1241$\,m）；本方案按要求采用均匀正七边形——覆盖与路径长度都不依赖朝向，
旋转自由度的几何意义更直观。

\begin{table}[H]
  \centering
  \caption{7 个观测站坐标（相对区域圆心）与半径：均匀正七边形布局，站点全部在 $r=1000$\,m 圆上}
  \label{tab:centers}
  \begin{tabular}{ccccc}
    \toprule
    站号 & $x$/m & $y$/m & 半径/m & 方位角/° \\
    \midrule
    0 & $+1000.00$ & $0.00$ & 1000.0 & 0.0 \\
    1 & $+623.49$ & $+781.83$ & 1000.0 & 51.4 \\
    2 & $-222.52$ & $+974.93$ & 1000.0 & 102.9 \\
    3 & $-900.97$ & $+433.88$ & 1000.0 & 154.3 \\
    4 & $-900.97$ & $-433.88$ & 1000.0 & 205.7 \\
    5 & $-222.52$ & $-974.93$ & 1000.0 & 257.1 \\
    6 & $+623.49$ & $-781.83$ & 1000.0 & 308.6 \\
    \bottomrule
  \end{tabular}
\end{table}

\subsection{巡视顺序与旋转自由度}
\begin{itemize}
  \item \textbf{巡视顺序}：以原点为起点的最短开放路径（Held--Karp 精确枚举 $7!$），长度
        $L_{\text{survey}}=6206.6$\,m，顺序 $2\to1\to0\to6\to5\to4\to3$（沿方位角单调扫过）；
  \item \textbf{旋转自由度}：把整个布局绕原点整体旋转，不改变覆盖保证（只依赖点间距离与点到
        原点距离），也不改变巡回路径长度（只依赖点间距离），故"各站朝向哪"是纯零成本自由度，
        用于把布局对准源密集方向（见下一节）。
\end{itemize}

%======================================================================
\section{起始全频道扫描与布局定向}
%======================================================================
\subsection{起点全频道扫描}
机器狗在原点（$0,0$）对全部 20 个频道各测向一次（里程 0，虚拟耗时 $\approx 20\times5+19\times1
=119$\,s），一次性拿到"哪些频道在接收范围内、从原点看它们各在哪个方向"。实测平均能锚定
$6\!\sim\!8$ 个源的方向，同时把其余频道标记为"源在 1000\,m 之外"（对后续跳过测量与区域收缩都有用）。

\subsection{最密集方向的估计}
设起始扫描听到的示向度集合为 $\mathcal{B}=\{b_1,\dots,b_m\}$。定义\textbf{最密集方向}为：
把圆周切成宽度 $60^\circ$ 的滑动窗，取落入窗口内示向度\textbf{个数最多}的窗口的\textbf{中心}方向
$\theta^*$（同分取角度最小者，确定性）：
\begin{equation}
  \theta^* = \arg\max_{\phi\in[0^\circ,360^\circ)}
  \bigl|\{b\in\mathcal{B} : \mathrm{ang\_diff}(b,\phi)\le 30^\circ\}\bigr|
  \label{eq:dense}
\end{equation}
其中 $\mathrm{ang\_diff}$ 取 $[0,180]$ 内的角度差。窗口逐 $1^\circ$ 扫描，随后在"装入个数不减少"
的前提下用窗口内示向度的圆均值做亚度精修。取\textbf{众数}而非均值：面对两簇相距几十度的等量源，
均值会落在两簇之间谁也没对准，众数才反映"源最多的一侧"。

\subsection{布局定向：第一点对准 $\theta^*-90^\circ$}
方案要求（用户指定）：把布局旋转，使\textbf{第一个巡视点}（最短开放路径的起点站）正对
"最密集方向\textbf{顺时针旋转 $90^\circ$}"的方位。在 atan2 坐标系（$x$ 右、$y$ 上、逆时针为正）
中，顺时针 $90^\circ$ 即方位角\textbf{减去 $90^\circ$}，故目标方位
\begin{equation}
  \theta_{\text{first}} = \theta^* - 90^\circ \pmod{360^\circ} .
  \label{eq:target}
\end{equation}
若第一点站原方位为 $\phi_0$，则布局整体旋转
\begin{equation}
  \Delta\varphi = \theta_{\text{first}} - \phi_0 \pmod{360^\circ},
  \label{eq:rotate}
\end{equation}
旋转后第一个巡视点正对 $\theta_{\text{first}}$。设计意图：把起点摆到源密集方向的旁侧 $90^\circ$，
使阶段一先扫外围开阔方向、把源密簇留给后续站位与阶段二（覆盖与路径长度均不受旋转影响，
纯零成本）。

%======================================================================
\section{阶段一：巡视扫描与顺路清除}
%======================================================================
\subsection{逐站测向与"够准即停"裁剪}
对每个巡视站，定义\textbf{活跃频道集}：
\begin{equation}
  \mathcal{A} = \bigl\{c : c\notin\text{cleared},\ |\text{obs}(c)|<3,\ c\text{ 不够准}\bigr\},
  \label{eq:active}
\end{equation}
即：未清除、示向度条数未达上限 $3$、且定位区域尚不够准的频道才测。\textbf{"够准"}定义为定位
区域的最小覆盖圆半径 $<20$\,m（等于清除半径）：叠加"无信号禁区 + 接收半径环带"后，很多频道
两条射线就把可能源集合压到 $<20$\,m，此时再测纯属浪费（每站 5\,s 测向 + 1\,s 切换）——实测
正是靠这条裁剪，使每站扫描量从第 1 站的 20 个频道逐站递减到 11 个。

每站对活跃频道按\textbf{编号升序}测向（升序省频道切换时间，最多省 $|\mathcal{A}|-1$ 次切换）；
测到 \texttt{near}（距离 $\le 5$\,m）就地清除；\texttt{no\_signal} 说明该站收不到该源，
会把"以站为圆心、以站到圆心连线方向为法线的半平面"并入该源定位区域的\textbf{无信号禁区}。

\subsection{顺路清除判据（方位扇区 + 双向）}
巡视途中去下一个站之前，利用"估计点"顺路清除。判据以区域圆心（原点）为参照：
\begin{itemize}
  \item \textbf{方位扇区}：估计点与圆心的连线方向 $b_{\text{est}}$ 位于"本站→圆心"方位
        $b_{\text{at}}$ 与"下一站→圆心"方位 $b_{\text{next}}$ 夹成的\textbf{较短弧}上，即
        \begin{equation}
          \mathrm{ang\_diff}(b_{\text{est}},b_{\text{at}})+\mathrm{ang\_diff}
          (b_{\text{est}},b_{\text{next}}) \;\le\;
          \mathrm{ang\_diff}(b_{\text{at}},b_{\text{next}})+\epsilon .
          \label{eq:arc}
        \end{equation}
        必须用角度而非 $\sin$（$|\sin 180^\circ|=0$ 会把反侧点误判为同向）；
  \item \textbf{半径上界}：估计点到圆心距离 $\le \max(\rho_{\text{at}},\rho_{\text{next}})$
        （"两点之间"自然止于两站所在半径，不延伸进圆域深处）；
  \item \textbf{600\,m 排除}：估计点距区域圆心 $<600$\,m（多为把原点当近似估计的不可靠点）不参与；
  \item \textbf{双向}：每站去下一站前，先后向清"当前站→上一站"之间（第 $k$ 站对第 $k-1$
        扇区补漏），再前向清"当前站→下一站"之间——两端各扫一遍，第一遍信息少可能漏，第二遍
        站在另一端估计更新后能补漏；首站无后向（起点段兜底覆盖），末站无前向；
  \item \textbf{起点段兜底}：从原点出发去第一站前，清"原点→第一站"方向 $\pm2^\circ$ 扇形内的
        估计点（原点是圆心，方位未定义，故退化为固定小扇形）。
\end{itemize}
清一个点要走到估计点再回来，成本固定（命中 5\,s / 未命中 3\,s 的清除 + 往返里程）；
命中后该频道后续各站都不再测量，省下阶段二专程跑的整段折返。

%======================================================================
\section{阶段二：定位与清除}
%======================================================================
对剩余每个频道，按"以估计点为中心的精确最短开放路径"顺序（Held--Karp 精确解，$n\le16$ 时
毫秒级）逐个处理。每个频道走四级流程：

\subsection{定位区域}
对每个频道，可能源集合 = 各观测示向度的 ${\pm1^\circ}$ 楔形之交 $\cap$ 圆域 $\cap$
无信号禁区之外；用内接 256 边形逼近圆域求交（误差 $<2$\,mm），取交区域的最小覆盖圆圆心作为
\textbf{估计点}。

\subsection{四级清除流程}
\begin{enumerate}
  \item \textbf{就近试清}：最小覆盖圆半径 $\le 400$\,m（估计不离谱，值得跑过去）时，走到估计点
        \code{/clear} 一次。实测 $90\%$ 的源在阶段二\textbf{一次就近试清即命中}；
  \item \textbf{多点覆盖试清}：未命中则用 $K=4$ 个半径 $20$\,m 的圆覆盖整个定位区域，依次试清
        （覆盖区域而非仅指望圆心，容错定位偏差）；
  \item \textbf{文献准则补测}：仍不行则按文献准则补测缩小区域：沿 $24$ 个方位角 $\times$
        $5$ 档半径（$150,300,450,600,800$\,m）构造候选补测点，优先选"预测 $\sigma$ 最小"的
        方向；一轮最多试 $3$ 个候选（收不到信号就换），最多 $6$ 轮；补测点与已有检测点保持
        $\ge 60$\,m 间距（同点复测无新信息）；
  \item \textbf{末端逼近}：兜底——沿最新实测示向度以 $16$\,m 步长逼近，步数按"当前位置到定位
        区域最远顶点距离"自适应（最少 24 步、上限 120 步），消除沿视线方向的定位偏差，
        保证"定位略偏也能清除"。
\end{enumerate}

%======================================================================
\section{确定性保证与验证}
%======================================================================
\subsection{确定性}
方案中所有决策（站点布局、巡视顺序、密集方向、旋转角、顺路清除判定、阶段二访问顺序、补测
候选、逼近步数）全部确定性：无随机数、并列取小编号规则固定、Held--Karp 遍历顺序固定。因此
\textbf{同一随机种子下整局逐字节可复现}——这是 A/B 对比与论文可复现性的基础。

\subsection{验证}
确定性方案不依赖独立脚本套件，采用三层可复跑检验：
\begin{itemize}
  \item \textbf{覆盖自检}：\code{python -m cumcm.t3.covering} 对默认 7 点布局与经典六边形族在
        100 个旋转角下复核最坏最近距离（与解析值一致、未破坏 1000\,m 保证），并复核密集扇区
        选向与"最少 7 个圆"的判定（6 个被已证明的 $\rho_6$ 排除），全部通过；
  \item \textbf{演练真值核对}：\code{--practice} 模式由 jammers-py 自带真值，逐局核对清除数、
        定位误差与 20\,m 清除半径 —— 20 局（seed 0--19）256/256 全清、最坏定位误差 19.24\,m
        $<20$\,m（见 \tabref{tab:main}）；
  \item \textbf{路线最优性}：巡视顺序由 Held--Karp 精确枚举最短开放路径（7 点精确），
        最近邻 + 2-opt 精修不使路径变长。
\end{itemize}
同一 seed 下全部决策逐字节可复现（覆盖求解、巡视顺序、顺路清除判定、补测候选与逼近步数
均无随机数）。

%======================================================================
\section{实测结果}
%======================================================================
\subsection{20 局统计（seed 0--19，离线演练同源模拟器）}
\begin{table}[H]
  \centering
  \caption{20 局（$N=256$ 个干扰源）总体成绩}
  \label{tab:main}
  \begin{tabular}{lcc}
    \toprule
    指标 & 数值 & 说明 \\
    \midrule
    清除比例 & $256/256=100\%$ & 全部命中 \\
    虚拟总时长 & 3590\,s & 范围 2558--4312\,s \\
    里程 & 13980\,m & 阶段一 6.2\,km + 阶段二 \\
    测向次数 & 121 次/局 & 够准即停裁剪后 \\
    定位误差 & 均值 7.69\,m / 最差 19.24\,m & $<$ 清除半径 20\,m \\
    顺路清除 & 命中 110 / 白跑 25 & 顺路命中率高、白跑少 \\
    阶段二处理 & 93\% 一次命中（238/256）& 补测后再命中 15、多源交会 3 \\
    \bottomrule
  \end{tabular}
\end{table}

\subsection{虚拟时长构成}
以 $5$\,m/s 移速折算，虚拟总时长 $3590$\,s 的构成：
\begin{itemize}
  \item \textbf{移动：约 78\%}（$2796$\,s）——大头，路径已双最优（巡视 Held--Karp、阶段二
        精确最短开放路径），逼近下界；
  \item \textbf{测向：约 17\%}（$606$\,s）——已被"够准即停""无信号禁区跳过""升序省切换"充分裁剪；
  \item 切换/清除/规划：约 5\%。
\end{itemize}

\subsection{与历史变体的对比（旧参数下的 10 局 seed 0--9，仅作策略演进参考）}
\begin{table}[H]
  \centering
  \caption{关键策略变体对比（同场景、同 seed，逐字节可复现；\textbf{绝对值为 2026-09-12 参数
  修订之前}，当前定案的 20 局数据见表~\ref{tab:main}）}
  \label{tab:variants}
  \begin{tabular}{lcccc}
    \toprule
    版本 & 平均虚拟/s & 顺路清 命中/白跑 & 备注 \\
    \midrule
    半径 1800 + 补测（全补） & 4473 & 89/58 & 补测 641 次，过重 \\
    MEC $\le 60$ 门槛 & 4057 & 85/12 & 白跑最少但命中略降 \\
    半径 1\,km 双向（旧朝向） & 3727 & 21/2 & 终点正对密集方向 \\
    半径 1\,km + 第一点 $\theta^*-90^\circ$ & \textbf{3688} & 20/4 & 旧参数下的本方案 \\
    \bottomrule
  \end{tabular}
\end{table}

%======================================================================
\section{关键参数}
%======================================================================
\begin{table}[H]
  \centering
  \caption{核心参数表}
  \label{tab:params}
  \begin{tabular}{lll}
    \toprule
    参数 & 值 & 含义 \\
    \midrule
    $\rho_{\text{region}}$ & 1800\,m & 目标圆域半径 \\
    $\rho_{\text{cover}}$ & 1000\,m & 覆盖圆半径 = 接收半径下界 \\
    $\rho_{\text{rec,max}}$ & 1500\,m & 接收半径上界 \\
    $r_{\text{clear}}$ & 20\,m & 清除半径 \\
    $|e|_{\max}$ & $1^\circ$ & 示向度误差上限 \\
    $v$ & 5\,m/s & 移动速度 \\
    $N_{\text{obs}}$ & $3$ & 单频道示向度采集上限 \\
    $60^\circ$ & —— & 密集方向窗口宽度 \\
    $\theta^* - 90^\circ$ & —— & 第一点目标方位（顺时针 90°） \\
    $r_{\text{excl}}$ & 600\,m & 顺路清除近心排除 \\
    $r_{\text{try}}$ & 400\,m & 就近试清的最大覆盖圆半径 \\
    $K_{\text{clear}}$ & 4 & 多点覆盖试清个数 \\
    $d_{\text{probe}}$ & 150--800\,m（5 档） & 补测候选半径 \\
    $h_{\text{step}}/h_{\max}/h_{\text{cap}}$ & 16\,m / 24 / 120 & 末端逼近步长/最少/上限 \\
    \bottomrule
  \end{tabular}
\end{table}

%======================================================================
\section{结论}
%======================================================================
本方案用"确定性两阶段 + 零成本布局定向"平衡了信息获取与移动成本：
覆盖布局保证每一源都能被听到；起点全频道扫描 + 密集方向定向把巡视摆到源密集方向旁侧；
逐站"够准即停"裁剪测向；顺路清除把可判定的源在巡视途中带走；阶段二四级流程（就近试清 → 
多点覆盖 → 文献补测 → 末端逼近）在 90\% 一枪中命中的同时，靠末端逼近兜底保证"定位略偏也能
清除"。

实测 20 局 256/256 全部清除，虚拟时长 3590\,s（移动占 78\% 且路径已逼近下界），最坏定位误差
19.24\,m $<$ 20\,m，覆盖自检与演练真值核对全部通过、同种子逐字节可复现。$\square$

\end{document}